\documentclass[../DoAn.tex]{subfiles}
\begin{document}

\begin{center}
    \Large{\textbf{TÓM TẮT NỘI DUNG ĐỒ ÁN}}\\
\end{center}
\vspace{1cm}
Trong bối cảnh rủi ro an ninh mạng ngày càng phức tạp, việc bảo vệ các hệ thống liên lạc nội bộ đang trở thành ưu tiên hàng đầu, thúc đẩy sự phổ biến của tiêu chuẩn mã hóa đầu cuối (E2EE) như một giải pháp bảo mật tất yếu. Hiện nay, hầu hết các ứng dụng nhắn tin E2EE hàng đầu vận hành dựa trên giao thức Double Ratchet, một tiêu chuẩn bộc lộ nhiều hạn chế khi mở rộng quy mô nhóm. Nhằm giải quyết bài toán này, Messaging Layer Security (MLS) ra đời như một giao thức tiên phong, tái định nghĩa toàn diện kiến trúc mã hóa nhóm. Tuy nhiên, tiêu chuẩn MLS hiện tại được thiết kế độc quyền cho kiến trúc tập trung và bắt buộc phụ thuộc vào một máy chủ điều phối để đảm bảo thứ tự. Việc triển khai MLS trên môi trường mạng ngang hàng (P2P) phi tập trung vẫn là một thách thức chưa được giải quyết. 

Từ thực tiễn đó, em quyết định theo đuổi hướng nghiên cứu tích hợp trực tiếp giao thức MLS lên kiến trúc mạng ngang hàng, tước bỏ hoàn toàn máy chủ trung tâm. Hướng đi này được lựa chọn nhằm chứng minh khả năng dung hợp giữa tiêu chuẩn mã hóa nhóm tiên tiến nhất hiện nay và triết lý mạng phi tập trung, từ đó vạch ra mô hình lý thuyết cho một hệ thống truyền thông không tồn tại điểm lỗi duy nhất. 

Tổng quan giải pháp của nghiên cứu là đề xuất một lớp điều phối phi tập trung. Kiến trúc này thay thế hoàn toàn vai trò của máy chủ trong việc quản lý cập nhật trạng thái nhóm, đồng thời cung cấp khả năng tự động duy trì trạng thái mã hóa nhóm, nhận diện và Fork Healing khi mạng xảy ra đứt gãy. Bên cạnh phần thiết kế giao thức, đồ án còn xây dựng một ứng dụng desktop nhằm hiện thực các luồng khởi tạo định danh, cấp bundle, tham gia tổ chức, trao đổi thông điệp mã hóa đầu cuối và quản trị thành viên. Sự hiện diện của ứng dụng này cho phép đối chiếu trực tiếp giữa mô hình giao thức và hành vi vận hành ở mức ứng dụng, qua đó củng cố tính khả thi của hướng tiếp cận được đề xuất.

\begin{flushright}
Sinh viên thực hiện\\
\vspace{1.5cm}
\begin{tabular}{@{}c@{}}
\textbf{Lê Tiến Dũng}
\end{tabular}
\end{flushright}

\end{document}
